\section{Plano de ação resumido}
\label{sec:plano}

A Tabela~\ref{tab:plano} resume o plano nos quatro elementos exigidos pela orientação: estratégia, ações-chave, recursos necessários e indicadores de sucesso. As metas são hipóteses derivadas da Atividade IX e deverão ser recalibradas após o piloto. Os responsáveis seguem a estrutura organizacional da Atividade XI \cite{grupo2026a11}.

\begin{table}[htb]
\centering
\caption{Quadro-resumo do plano de ação}
\label{tab:plano}
\fontsizetable
\begin{tabularx}{\textwidth}{L{2.2cm} L{4.1cm} L{3.6cm} Y}
\toprule
\textbf{Estratégia} & \textbf{Ações-chave} & \textbf{Recursos necessários} & \textbf{Indicadores de sucesso (metas-hipótese)}\\
\midrule
\textbf{E1 -- Densidade local por parcerias (principal)} &
1. Fechar piloto com 3 clínicas (3 meses gratuitos do Plano Parceiro) e 2 ONGs ou lideranças comunitárias (Giovanna).\newline
2. Instalar código QR no balcão e iniciar divulgação em redes sociais e indicação com créditos (Júlia).\newline
3. Lançar piloto com tutores reais após correções de segurança e termos de privacidade (Gabriel, Kathleen, Júlia).\newline
4. Medir o CAC e a retenção; aplicar os gatilhos da Tabela~\ref{tab:gatilhos} antes de abrir nova praça. &
Tempo dos quatro integrantes; materiais de divulgação (R\$~150,00 mensais já previstos) e teto de aquisição (cerca de R\$~438 mensais); consultoria veterinária (R\$~600,00 mensais) e jurídico e LGPD (R\$~250,00 mensais), já nos custos fixos; infraestrutura em nuvem (R\$~100,00 mensais). &
Usuários ativos: 300 (mês 3), 800 (mês 6), 1.500 (mês 12).\newline
Clínicas ativas: 3, 6 e 10 nesses mesmos marcos.\newline
Crescimento mensal de usuários: cerca de 20\% (meses 3 a 12).\newline
CAC observado $\le$ teto (R\$~3,29).\newline
Atendimentos intermediados: 10\% dos usuários ativos por mês.\\
\midrule
\textbf{E2 -- Aprofundamento de receita e novos serviços (complementar)} &
1. Aplicar o teste de aceitação do preço da Atividade IX (Júlia, Giovanna).\newline
2. Lançar o PetCuida+ com plano anual e meta de conversão de 5\% a 8\% (Giovanna, Gabriel).\newline
3. Vender campanhas patrocinadas e ao menos um convênio para financiar o fundo (Giovanna).\newline
4. Priorizar novos serviços da lista de espera somente com demanda comprovada (decisão do colegiado). &
Tempo de desenvolvimento (Gabriel e Kathleen); taxa de conta de desenvolvedor da loja de aplicativos; material comercial para patrocinadores; parecer jurídico para contratos e termos. &
Conversão ao PetCuida+: $\ge$ 5\% (meta 8\%).\newline
Fundo de liquidação por atendimento: de R\$~6,60 para R\$~3,30.\newline
Receita recorrente: $\ge$ 40\% da receita (sem convênios).\newline
Convênio ou patrocínio contratado: $\ge$ 1.\newline
Contribuição por usuário: de R\$~0,55 para R\$~1,42.\\
\bottomrule
\end{tabularx}
\fonteautores
\end{table}

\subsection{Cronograma e marcos}

\begin{alineas}
  \item \textbf{Meses 0 a 3 (preparação e piloto):} correções de segurança e privacidade; constituição da empresa; teste de aceitação do preço; três clínicas e duas ONGs; meta de 300 usuários ativos;
  \item \textbf{meses 4 a 6 (validação do núcleo):} medição de CAC, retenção e atendimentos por usuário; seis clínicas e 800 usuários; revisão do fundo de liquidação;
  \item \textbf{meses 7 a 12 (consolidação):} dez clínicas, 1.500 usuários, 75 assinantes e ao menos um convênio; decisão colegiada sobre abertura de nova praça;
  \item \textbf{meses 13 a 24 (expansão condicionada):} abertura de uma nova área da região metropolitana se os gatilhos forem atingidos; meta de 5.000 usuários no mês 24; avaliação de financiamento externo, pois o cenário com equipe remunerada depende de escala maior ou de recursos de terceiros.
\end{alineas}

\subsection{Governança do plano}

A execução segue as regras da Atividade XI: Giovanna aprova escopo, custos e parcerias; Gabriel, a arquitetura; Júlia, privacidade, interface e comunicação externa; e Kathleen, a liberação de versões e a facilitação do alinhamento semanal. A revisão mensal de negócio e de carga compara os indicadores da Tabela~\ref{tab:plano} com as metas e decide se o plano deve ser mantido, ajustado ou interrompido.
